\section{Common charts for moving return domains}
\label{sec:common-charts}
We give a local change-of-variables estimate and then apply it to the
actual return domains. This makes explicit the domination used in
Theorem~\ref{thm:return-stability}. Critical neighborhoods are discarded
in the initial collision space, not by assuming a uniform density bound
near their images.

\begin{lemma}[A common submersion chart]
\label{lem:common-submersion-chart}
Let $U=U_1\times U_2\Subset\R^2$ be a bounded open rectangle and let
$f(R,u,v)$ be $C^2$ on a neighborhood of $J\times\overline U$, where
$J$ is a compact parameter interval. Suppose
$\partial_v f\ge\eta>0$ there. For $a\in C_c^1(U)$, let $g_R$ be the
density of the pushforward of $a(u,v)\dd u\dd v$ by $f_R$. Then
\begin{equation}\label{eq:chart-Lipschitz}
 \|g_R-g_S\|_{L^1(\R)}\le C_a|R-S|\qquad(R,S\in J).
\end{equation}
The constant depends only on $U$, a common bounded range of $f$, $\eta$,
the $C^2$ bound for $f$, and the $C^1$ bound and support of $a$.
The same assertion holds if $\partial_v f\le-\eta$.
\end{lemma}
\begin{proof}
For fixed $R,u$ the map $v\mapsto f(R,u,v)$ is one-to-one. Write its
inverse as $v_R(u,t)$. The pushforward under $(u,v)\mapsto(u,f_R(u,v))$
has density
\[
 A_R(u,t)=\frac{a(u,v_R(u,t))}{\partial_v f(R,u,v_R(u,t))}
\]
on its image, extended by zero elsewhere. All images lie in one bounded
$(u,t)$ box. Since $a$ vanishes on a fixed neighborhood of $\partial U$,
this zero extension introduces no boundary jump as $R$ varies. Wherever
$a$ or its derivatives are nonzero, the inverse is defined on a common
parameter neighborhood. If $M$ bounds the relevant derivatives of $f$,
implicit differentiation gives
\[
 |\partial_R v_R|\le M/\eta,\qquad
 |\partial_R A_R|
 \le \|a_v\|_\infty M/\eta^2
       +\|a\|_\infty(M/\eta^2+M^2/\eta^3).
\]
At points outside the inverse chart both sides are interpreted by the
zero extension, which vanishes on an open neighborhood of its moving
edge. The same bound therefore gives an $L^1$ majorant on the fixed box.
Integrating in $R$ proves
$\|A_R-A_S\|_1\le C_a|R-S|$. Since
$g_R(t)=\int A_R(u,t)\dd u$, integration in $u$ proves
\eqref{eq:chart-Lipschitz}. Reversing the $v$ coordinate handles the
negative derivative. The result also applies on a constant lattice
component of a mixed density.
\end{proof}

\begin{proposition}[Finite common patches with an explicit remainder]
\label{prop:finite-common-patches}
Fix $n\ge1$, $R_0\in I$, and $\epsilon>0$. There are finitely many
nonnegative smooth cutoffs $\psi_i$ on the common collision space, a
parameter neighborhood $J$ of $R_0$, and $C_{\epsilon,n,R_0}<\infty$
such that $\sum_i\psi_i\le1$, every support lies in one common regular
first-$n$-return patch for all $R\in J$, and
\begin{equation}\label{eq:common-patch-mass}
 c_*^{-1}\int\sum_i\psi_i\dd\nu>1-\epsilon.
\end{equation}
The lattice labels and all intermediate return decisions are fixed on
each patch. If $p^i_{n,R}$ is the pushforward density of
$c_*^{-1}\psi_i\nu$, then
\begin{align}\label{eq:common-patch-bound}
 \sum_i\|p^i_{n,R}-p^i_{n,S}\|_1
       &\le C_{\epsilon,n,R_0}|R-S|,\\
 \|p_{n,R}-p_{n,S}\|_1
       &\le2\epsilon+C_{\epsilon,n,R_0}|R-S|
               \qquad(R,S\in J).\nonumber
\end{align}
Before choosing these patches one may discard $N_{n,R_0}>L$ at cost at
most $C_1e^{-a_1L/n}$, with the constants of
Theorem~\ref{thm:exponential-return}.
\end{proposition}
\begin{proof}
The exponential moment bound proves the last assertion. Choose $L$ so
that this cost is less than $\epsilon/3$. On the remainder remove the
initial boundary of the section, all grazing and competing-root
singularities encountered before collision $L$, and the preimages of the
moving return boundary at $R_0$ through those collisions. These sets have
zero collision measure. On every remaining regular word also remove the
critical points of its cumulative flight length. Proposition
\ref{prop:general-edge} shows that there is at most one such point on a
regular optical word; the submersion complement has full measure.
There are only finitely many possible physical lattice words of length
at most $L$, since each flight label belongs to a fixed finite set.
No estimate for the number of unbounded-length words is used here.

Inner regularity of the collision measure and a countable exhaustion of
the regular submersion patches give a compact subset of mass greater
than $1-\epsilon$ in the normalized section. At every point of this
compact subset, choose a small coordinate rectangle on which one
coordinate derivative of the cumulative time is bounded away from zero.
The finite-itinerary lemma preserves the entire physical word and every
return decision on a slightly larger rectangle for $R$ close to $R_0$.
Select finitely many such rectangles and a subordinate family of smooth
cutoffs with sum at most one, equal to one on the compact subset.
Their supports have positive distance from all the excluded sets at
$R_0$. Finiteness gives a common parameter neighborhood $J$ on which
these distances, return decisions and derivative bounds persist. The
time functions and their first two derivatives are uniformly bounded
on these finitely many compact rectangles. Lemma
\ref{lem:common-submersion-chart} applies with
$a=\psi_i/(4\pi c_*)$ in $(\alpha,p)$ coordinates. Summing its bounds
proves the first line of \eqref{eq:common-patch-bound}.

For every $R\in J$ the remainder of the original probability is a
positive measure. Its mass is exactly
$1-c_*^{-1}\int\sum_i\psi_i\dd\nu<\epsilon$, because the section
measure is the fixed number $c_*$. The raw densities exist by the
submersion/coarea argument of Theorem~\ref{thm:return-stability}.
The two remainder densities have total $L^1$ norm less than
$2\epsilon$. This gives the second line. Nothing has been assumed about
the size of a density at a critical value or a sharp image boundary.
\end{proof}

The constants in \eqref{eq:common-patch-bound} expose the order of
choices: physical-count truncation, deletion of small initial-state
neighborhoods, compact common charts, and finally parameter tolerance.
They are not asserted to remain bounded as $n$ or $\epsilon^{-1}$ grows.
A quantitative global density modulus requires quantitative control of
those additional choices. Bounded measurable insertions are handled by
$L^1$ approximation on the same patches, as in
Proposition~\ref{prop:weighted-return-stability}; no artificial observation
coordinate is introduced.
